\documentclass[11pt]{article}
\usepackage{preamble}
\usepackage{multicol}
\setlength{\parskip}{4pt}

\begin{document}

\daybanner{AP Statistics \textperiodcentered\ Unit 1 \textperiodcentered\ Topics 1.7--1.8 (CED 1.7--1.8) \textperiodcentered\ TEACHER KEY}{Typical Hours, Exceptional Week}

\sectionbanner{CED TETHER}
\begin{multicols}{2}
{\footnotesize\setlength{\parskip}{0pt}
\begin{itemize}[leftmargin=*,itemsep=0pt,topsep=0pt]
\item ENDURING UNDERSTANDING: UNC-1 Graphical representations and statistics allow us to identify and represent key features of data.
\item SKILLS:
\item \textbf{Skill 2.C} --- Calculate summary statistics, relative positions of points within a distribution, correlation, and predicted response.
\item \textbf{Skill 4.B} --- Interpret statistical calculations and findings to assign meaning or assess a claim.
\item LO UNC-1.I Calculate measures of center and position for quantitative data. [Skill 2.C]
\item EK UNC-1.I.1 A statistic is a numerical summary of sample data.
\item EK UNC-1.I.2 The mean is the sum of all the data values divided by the number of values. For a sample, the mean is denoted by x-bar: $\bar{x}$ = (1/n) $\sum$ $x_i$, where $x_i$ represents the $i^{\mathrm{th}}$ data point in the sample and n represents the number of data values in the sample.
\item EK UNC-1.I.3 The median of a data set is the middle value when data are ordered. When the number of data points is even, the median can take on any value between the two middle values. In AP Statistics, the most commonly used value for the median of a data set with an even number of values is the average of the two middle values.
\item EK UNC-1.I.4 The first quartile, Q1, is the median of the half of the ordered data set from the minimum to the position of the median. The third quartile, Q3, is the median of the half of the ordered data set from the position of the median to the maximum. Q1 and Q3 form the boundaries for the middle 50\% of values in an ordered data set.
\item EK UNC-1.I.5 The p\textsuperscript{th} percentile is interpreted as the value that has p\% of the data less than or equal to it.
\item LO UNC-1.J Calculate measures of variability for quantitative data. [Skill 2.C]
\item EK UNC-1.J.1 Three commonly used measures of variability (or spread) in a distribution are the range, interquartile range, and standard deviation.
\item EK UNC-1.J.2 The range is defined as the difference between the maximum data value and the minimum data value. The interquartile range (IQR) is defined as the difference between the third and first quartiles: Q3 - Q1. Both the range and the interquartile range are possible ways of measuring variability of the distribution of a quantitative variable.
\item EK UNC-1.J.3 Standard deviation is a way to measure variability of the distribution of a quantitative variable. For a sample, the standard deviation is denoted by $s_x$: $s_x$ = $\surd$[(1/(n-1)) $\sum$($x_i$ - $\bar{x}$)\textsuperscript{2}]. The square of the sample standard deviation, $s^2$, is called the sample variance.
\item EK UNC-1.J.4 Changing units of measurement affects the values of the calculated statistics.
\item LO UNC-1.K Explain the selection of a particular measure of center and/or variability for describing a set of quantitative data. [Skill 4.B]
\item EK UNC-1.K.1 There are many methods for determining outliers. Two methods frequently used in this course are: i. An outlier is a value greater than 1.5 $\times$ IQR above the third quartile or more than 1.5 $\times$ IQR below the first quartile.
\item ii. An outlier is a value located 2 or more standard deviations above, or below, the mean.
\item EK UNC-1.K.2 The mean, standard deviation, and range are considered nonresistant (or non-robust) because they are influenced by outliers. The median and IQR are considered resistant (or robust), because outliers do not greatly (if at all) affect their value.
\item Enduring Understanding: UNC-1 Graphical representations and statistics allow us to identify and represent key features of data.
\item Skills:
\item \textbf{Skill 2.B} --- Construct numerical or graphical representations of distributions.
\item \textbf{Skill 2.A} --- Describe data presented numerically or graphically.
\item Learning Objectives and Essential Knowledge:
\item UNC-1.L Represent summary statistics for quantitative data graphically. [Skill 2.B]
\item \textbf{UNC-1.L.1} --- Taken together, the minimum data value, the first quartile (Q1), the median, the third quartile (Q3), and the maximum data value make up the five-number summary.
\item \textbf{UNC-1.L.2} --- A boxplot is a graphical representation of the five-number summary (minimum, first quartile, median, third quartile, maximum). The box represents the middle 50\% of data, with a line at the median and the ends of the box corresponding to the quartiles. Lines ("whiskers") extend from the quartiles to the most extreme point that is not an outlier, and outliers are indicated by their own symbol beyond this.
\item UNC-1.M Describe summary statistics of quantitative data represented graphically. [Skill 2.A]
\item \textbf{UNC-1.M.1} --- Summary statistics of quantitative data, or of sets of quantitative data, can be used to justify claims about the data in context.
\item \textbf{UNC-1.M.2} --- If a distribution is relatively symmetric, then the mean and median are relatively close to one another. If a distribution is skewed right, then the mean is usually to the right of the median. If the distribution is skewed left, then the mean is usually to the left of the median.
\end{itemize}}
\end{multicols}
\clearpage
\small
\textbf{Essential question:} Which numerical and graphical summary communicates a typical week without hiding an exception?\par
\textbf{DOK-3 skill:} 4.B \textperiodcentered\ \textbf{FRQ pattern:} {\small\texttt{adjudicate-\allowbreak{}competing-\allowbreak{}displays-\allowbreak{}and-\allowbreak{}bound-\allowbreak{}the-\allowbreak{}report}}\par
\textbf{DOK rationale:} Part (c) coordinates the two topics to judge competing claims, defend a corrected report, and explain what the rejected representation cannot establish.\par

\textbf{Self-paced bonus sheet.} Students take it when they choose and turn it in when they finish. Everything they need is printed on the sheet. Score part (c) only, E / P / I, by hand --- never AI-graded or auto-scored. (a) and (b) are the ladder up; use them to see where a thin (c) came from.\par

\textbf{Questions to ask}
\begin{itemize}[leftmargin=1.5em,itemsep=2pt,topsep=2pt]
  \item Which feature settles that claim?
  \item What information would your chosen display hide?
\end{itemize}

\begin{callout}{\IconWarn\ LOOK FOR}{calloutred}
\begin{itemize}[leftmargin=1.5em,itemsep=2pt,topsep=2pt]
  \item Extending the upper whisker to 28 instead of marking it separately.
  \item Equating mean with the value most students report.
  \item Claiming that a boxplot preserves each repeated observation.
\end{itemize}
\end{callout}

\sectionbanner{SCORING PART (c) --- E / P / I}

\textbf{Expected elements}
\begin{itemize}[leftmargin=1.5em,itemsep=2pt,topsep=2pt]
  \item \textbf{e1} (required) --- Chooses a defensible center/spread pair using 28 and resistance, distinguishing correct arithmetic from representativeness.
  \item \textbf{e2} (required) --- Corrects the claim that a boxplot retains every individual value.
  \item \textbf{e3} (required) --- Defends a display choice with a benefit and a specific information limit in context.
\end{itemize}

{\small\renewcommand{\arraystretch}{1.25}
\noindent\begin{tabular}{|p{0.5in}|p{5.9in}|}
\hline
\textbf{E} & All three evidence-to-reason links are correct in context. Accept an equivalent defensible report. \\ \hline
\textbf{P} & At least one correct evidence-to-reason link, but another is missing or incorrect. \\ \hline
\textbf{I} & No correct evidence-to-reason link; only a choice, copied numbers, or unsupported claims. \\ \hline
\end{tabular}}

\textbf{Common mistakes}
\begin{itemize}[leftmargin=1.5em,itemsep=2pt,topsep=2pt]
  \item Extending the upper whisker to 28 instead of marking it separately.
  \item Equating mean with the value most students report.
  \item Claiming that a boxplot preserves each repeated observation.
\end{itemize}

\clearpage
\normalsize
\focusbanner{THE PROBLEM --- annotated}

Suppose ten students report weekly paid-work hours: 2, 3, 4, 4, 5, 5, 6, 6, 7, 28. The mean is 7 hours.\par\smallskip \textbf{Zoe:} ``Use the median and IQR with a boxplot; the boxplot will keep every individual hour value. '' \textbf{Malik:} ``Report the mean: a typical student works 7 hours, because the mean uses every observation.''\par
\begin{center}
\begin{tikzpicture}
\begin{axis}[
  width=0.68\linewidth,
  height=1.28in,
  ymin=0, ymax=3, ytick=\empty, axis y line=none, axis x line=bottom,
  xlabel={Weekly paid-work hours}, tick label style={font=\small}, label style={font=\small}
]
\addplot[only marks, mark=*, mark size=1.9pt, color=dayblue] coordinates {
  (2,1)
  (3,1)
  (4,1)
  (4,2)
  (5,1)
  (5,2)
  (6,1)
  (6,2)
  (7,1)
  (28,1)
};
\end{axis}
\end{tikzpicture}
\end{center}
\begin{firsttakebox}{FIRST TAKE --- one sentence before you work the parts}
Which report would you send to the principal, or would you revise both? Cite a feature of the data.\par
\answer{Ungraded commitment; accept any evidence-based first impression.}
\end{firsttakebox}

\begin{callout}{\IconBook\ TWO TOPICS, ONE DATA STORY}{calloutgreen}
1.7: Median and IQR resist extremes; mean and SD are more sensitive. Here quartiles are the medians of the two halves.\\ 1.8: A modified boxplot marks values beyond $Q_1-1.5\mathrm{IQR}$ or $Q_3+1.5\mathrm{IQR}$ separately.\\ Whiskers reach the most extreme nonoutliers; the box spans the middle half and can hide repeats and gaps.
\end{callout}

\dokbadge{1}~\textbf{(a)} State the median weekly work hours.\par
\answer{The two middle values are both 5, so the median is 5 hours.}

\dokbadge{2}~\textbf{(b)} Use medians of the lower and upper five values to find $Q_1$, $Q_3$, and IQR. Apply the $1.5\,\mathrm{IQR}$ rule and sketch a modified boxplot, labeling the whiskers and any outlier.\par
\answer{$Q_1=4$, $Q_3=6$, IQR $=2$. Fences are $1$ and $9$ hours. The box runs 4 to 6 with median 5; whiskers reach 2 and 7; mark 28 separately as an outlier.}

\dokbadge{3}~\textbf{(c)} $\star$ In three to five sentences, evaluate both reports and propose a corrected report of typical hours and spread. Use the unusual value to explain your summary choice. Decide whether to keep the boxplot or dotplot, stating one benefit and one information limit of your choice.\par
\answer{Report median 5 and IQR 2 hours: the isolated 28-hour week pulls the mean upward to 7, while the resistant summaries describe the middle half. Zoe chooses a defensible summary pair but overstates the boxplot: it compresses the observations and loses repeats such as the two 4s, 5s, and 6s. Malik\textquotesingle s mean is correct but is less representative of the central cluster. A boxplot efficiently shows median, IQR and the outlier but loses exact frequencies; a dotplot retains those frequencies but requires reading or calculating the quartiles. Accept either display with its benefit and limit explained.}

\end{document}
